\documentclass[10pt]{article}
\usepackage[margin=1in]{geometry}
\usepackage{amsmath,amssymb,amsthm,mathtools}
\usepackage[hidelinks]{hyperref}
\newtheorem{theorem}{Theorem}
\newtheorem{lemma}[theorem]{Lemma}
\newtheorem{proposition}[theorem]{Proposition}
\newcommand{\E}{\mathbb E}
\newcommand{\PP}{\mathbb P}
\newcommand{\Be}{\operatorname{Beta}}
\title{Full real-shift log-concavity for gamma-ratio series\\
\large A proof of MathDB problem 324473}
\author{Principia Math In Progress}
\date{September 11, 2026}
\begin{document}
\maketitle

\noindent\textbf{Task 1: solve the original problem.}
We prove Conjecture 2 of Kalmykov--Karp \cite[p.~413]{KK},
with precisely the hypotheses of their Theorem 4. Here \emph{\(q\)-log-concave}
means that the coefficients of a product difference are nonnegative;
there is no basic-hypergeometric parameter \(q\).

\begin{theorem}\label{thm:main}
Let \(a\ge c>0\) and
\[
 G(\mu;x)=\sum_{j\ge0}g_j\frac{\Gamma(a+\mu+j)}{\Gamma(c+\mu+j)}
                    \frac{x^j}{j!}.
\]
Suppose either \(0\le a-c\le1\) and \(g_j\ge0\) are arbitrary, or
\(a-c>1\) and the nonnegative sequence \(g\) is log-concave without
internal zeros. Then, for every \(\mu,\nu\ge0\), every coefficient of
\[
 G(\mu;x)G(\nu;x)-G(0;x)G(\mu+\nu;x)
 \tag{1}
\]
is nonnegative. The statement holds for formal series, and consequently
also gives Wright log-concavity for nonnegative \(x\) in the common
domain of convergence.
\end{theorem}

We first prove a distributional statement that supplies the main case.
A probability law on \([0,1]\) is called symmetric if it is invariant under
\(z\mapsto1-z\). For symmetric laws write \(X\succeq Y\) if
\[
 \PP(a\le X\le1-a)\ge \PP(a\le Y\le1-a)\qquad(0\le a\le\tfrac12).
 \tag{2}
\]
By integrating superlevel sets, (2) implies
\(\E h(X)\ge\E h(Y)\) for every bounded nonnegative symmetric function
\(h\) that is nondecreasing on \([0,1/2]\). We call such functions
symmetrically unimodal. A transition kernel preserves (2) whenever it is
symmetric and the conditional probability of each centered interval is
symmetrically unimodal as a function of its starting point.

\begin{lemma}\label{lem:beta}
At fixed \(S=u+v>0\), the symmetrization of \(\Be(u,v)\) becomes more
concentrated in the order (2) as \(|u-v|\) decreases.
\end{lemma}
\begin{proof}
Put \(\delta=|u-v|/2<S/2\). The symmetrized density is
\[
 \frac{[z(1-z)]^{S/2-1}}{B(S/2-\delta,S/2+\delta)}
 \cosh\!\left(\delta\log\frac z{1-z}\right).
 \tag{3}
\]
For \(\delta_2\ge\delta_1\ge0\), the ratio of the two densities is
nondecreasing with \(|z-1/2|\). Indeed, for \(t\ge0\),
\[
 \frac{d}{dt}\log\frac{\cosh(\delta_2t)}{\cosh(\delta_1t)}
 =\delta_2\tanh(\delta_2t)-\delta_1\tanh(\delta_1t)\ge0.
\]
Both densities integrate to one, so this ratio can cross one only from
below to above as the distance from the center increases. Integration
over centered intervals proves the claim.
\end{proof}

\begin{lemma}\label{lem:coin}
For \(s>0\), let \(U\sim\Be(s,1)\) and let \(\epsilon\) be an independent
fair coin taking the values \(0,1\). The transition
\[
 x\longmapsto Ux+(1-U)\epsilon
 \tag{4}
\]
preserves (2).
\end{lemma}
\begin{proof}
The kernel is symmetric. For \(0<a<1/2\) and \(x\ge1/2\), its conditional
probability of \([a,1-a]\) is
\[
 p_s(x)=
 \begin{cases}
 1-\dfrac{a^s}{2}\bigl(x^{-s}+(1-x)^{-s}\bigr),
                          &\frac12\le x\le1-a,\\[4pt]
 \dfrac{(1-a)^s-a^s}{2x^s},&1-a\le x\le1.
 \end{cases}                                           \tag{5}
\]
This follows directly from \(\PP(U\le t)=t^s\), separately for each
coin outcome. Both pieces decrease with \(x\), and their values agree
at \(1-a\). Symmetry and the endpoint cases by limits prove the assertion.
\end{proof}

\begin{lemma}\label{lem:convolution}
Let \(Y\) have a symmetrically unimodal density on \([0,1]\). If \(U\)
is any independent random variable with values in \((0,1)\), then
\[
 x\longmapsto Ux+(1-U)Y
 \tag{6}
\]
preserves (2).
\end{lemma}
\begin{proof}
Fix \(U=\lambda\). Center all variables at \(1/2\). The probability that
\(\lambda(x-1/2)+(1-\lambda)(Y-1/2)\) belongs to a centered interval
is the mass of an interval of fixed length translated through a symmetric
density that decreases away from zero. That mass decreases with the
absolute value of its center: for a continuous density \(f\), its
derivative at a nonnegative center \(t\), for interval radius \(b\), is
\(f(t+b)-f(t-b)\le0\). Approximation gives the same conclusion for a
general unimodal density. Thus the conditional probability is
symmetrically unimodal in \(x\). Averaging over \(\lambda\) proves (6).
\end{proof}

We use elementary beta--gamma addition identities. If \(A,B\) are
independent gamma variables with respective positive shapes \(u,v\)
and unit rate, then \(A/(A+B)\sim\Be(u,v)\) is independent of
\(A+B\), whose shape is \(S=u+v\). These assertions follow by the change
of variables \(A=tz,\ B=t(1-z)\) in their product density. In particular,
for independent
\[
 X\sim\Be(u,v),\quad U\sim\Be(S,n),\quad
 Y\sim\Be(k,n-k),\qquad 1\le k\le n-1,
\]
we have
\[
 UX+(1-U)Y\sim\Be(u+k,v+n-k).                \tag{7}
\]
To see this, use independent gamma pairs of shapes \((u,v)\) and
\((k,n-k)\), then add corresponding components.

\begin{proposition}\label{prop:weights}
Fix \(n\ge0,\ d\ge0\), and nonnegative weights satisfying
\[
 w_k=w_{n-k},\qquad w_0\le w_1\le\cdots\le w_{\lfloor n/2\rfloor}.
 \tag{8}
\]
With \(r_d(z)=\Gamma(z+d)/\Gamma(z)\), the function
\[
 H_w(u,v)=\sum_{k=0}^n\binom nk w_k r_d(u+k)r_d(v+n-k),\qquad u,v>0,
 \tag{9}
\]
is nondecreasing as \(u,v\) are balanced at fixed sum.
\end{proposition}
\begin{proof}
The case \(d=0\) is constant in \(u,v\), so assume \(d>0\). Write
\(\phi_d(z)=[z(1-z)]^d\), a bounded symmetrically unimodal function.
For fixed \(S=u+v\), the beta integral gives
\[
 r_d(u+k)r_d(v+n-k)
 =\frac{\Gamma(S+n+2d)}{\Gamma(S+n)}
   \E\phi_d(B_{u+k,v+n-k}),                 \tag{10}
\]
where \(B_{p,q}\sim\Be(p,q)\). The multiplier depends only on \(S,n,d\).

First take \(w_k=1\). The mixture of beta laws with probabilities
\(2^{-n}\binom nk\) can be obtained by starting with independent
gamma variables of shapes \(u,v\) and adding \(n\) independent unit-rate
exponentials, each assigned to either component by a fair coin.
Conditional on \(k\) assignments to the first component, the normalized
first component has law \(\Be(u+k,v+n-k)\).
After \(j\) additions, the total has gamma shape \(S+j\) and is
independent of the normalized first component. Thus the next addition
acts on that component by (4) with \(s=S+j\), for \(j=0,\ldots,n-1\).
Symmetrizing the initial beta law does not change expectations of
symmetric functions after these symmetric transitions.
Lemma~\ref{lem:beta}, followed by \(n\) applications of
Lemma~\ref{lem:coin}, therefore proves the required monotonicity in
(10) for \(w_k=1\), including \(n=0\).

For general \(w\), put \(b_k=w_k-w_0\). Then
\[
 H_w=w_0H_{\boldsymbol 1}+H_b,\qquad b_0=b_n=0,
 \tag{11}
\]
and \(b\) still satisfies (8). If \(b\) vanishes identically there is
nothing more to prove; otherwise \(n\ge2\). Set
\[
 Z=\sum_{k=1}^{n-1}\binom nk b_k>0,\qquad
 \PP(K=k)=Z^{-1}\binom nk b_k.
\]
Conditional on \(K=k\), let \(Y\sim\Be(k,n-k)\). For the Bernstein
polynomials \(\mathcal B_{N,j}(y)=\binom Nj y^j(1-y)^{N-j}\), the density
of \(Y\) is
\[
 f_Y(y)=\frac{n-1}{Z}\sum_{j=0}^{n-2}\alpha_j\mathcal B_{n-2,j}(y),
 \qquad \alpha_j=\binom n{j+1}b_{j+1}.       \tag{12}
\]
The coefficients \(\alpha_j\) are symmetric and nondecreasing toward
their center, since both factors defining them have those properties.
Consequently (12) is symmetrically unimodal. For completeness, if
\(N=n-2\ge1\), differentiate the Bernstein sum to obtain
\[
 N\sum_{j=0}^{N-1}(\alpha_{j+1}-\alpha_j)\mathcal B_{N-1,j}(y).
 \tag{13}
\]
The coefficient differences are antisymmetric and nonnegative on
their lower-index half. Pair reflected indices in (13); for
\(y\le1/2\) the lower-index Bernstein polynomial is at least the
higher-index one. The derivative is therefore nonnegative. If \(N=0\),
the density is constant.

Take independent \(X\sim\Be(u,v)\), \(U\sim\Be(S,n)\), and this \(Y\).
Identity (7), conditional on \(K\), shows that \(UX+(1-U)Y\) is exactly
the beta mixture corresponding to \(H_b/Z\) in (10).
Its symmetrization becomes more concentrated under balancing, by
Lemmas~\ref{lem:beta} and \ref{lem:convolution}; the law of \(U\) depends
only on the fixed sum \(S\). Its expectation of \(\phi_d\) thus increases.
This proves the assertion for \(H_b\), and (11) completes the proof.
\end{proof}

\begin{proof}[Proof of Theorem~\ref{thm:main}]
Put \(d=a-c\). If \(g\) is log-concave without internal zeros, the
weights \(w_k=g_kg_{n-k}\) satisfy (8). Indeed, on their positive support,
\[
 \frac{w_{k+1}}{w_k}
 =\frac{g_{k+1}/g_k}{g_{n-k}/g_{n-k-1}}\ge1\quad(k<n/2),
 \tag{14}
\]
because successive ratios of a log-concave sequence decrease.
The positive support of \(w\) is a centered interval, so zeros cause
no exception; identically zero weights are harmless.

By the Cauchy product, \(n!\) times the coefficient of \(x^n\) in (1) is
\[
 H_w(c+\mu,c+\nu)-H_w(c,c+\mu+\nu).           \tag{15}
\]
The pairs have equal sums, and the first has smaller absolute
difference since \(|\mu-\nu|\le\mu+\nu\). Proposition~\ref{prop:weights}
proves (15) nonnegative. This proves the required \(d>1\) case, and
also proves the result for log-concave \(g\) for every \(d\ge0\).

It remains to allow arbitrary nonnegative \(g\) when \(0\le d\le1\).
For \(0<d<1\), integration by parts and the beta integral give
\[
 r_d(z)=\frac d{\Gamma(1-d)}
 \int_0^1 (1-t^z)t^{d-1}(1-t)^{-d-1}\,dt.     \tag{16}
\]
The boundary terms vanish: at zero the relevant factor is \(t^d\),
and at one it is \(O((1-t)^{1-d})\). Differentiating under this
convergent integral shows \(r_d'(z)>0\) and \(r_d''(z)<0\).
For \(d=1\), \(r_d(z)=z\) is increasing and concave as well.

If \(r\) is nonnegative, increasing, differentiable and concave, then,
for any \(k,l\ge0\), its symmetrized product
\[
 P(u,v)=r(u+k)r(v+l)+r(v+k)r(u+l)
\]
increases under balancing. For \(u<v\), group its directional derivative as
\begin{align*}
 (\partial_u-\partial_v)P
 &=\bigl[r'(u+k)r(v+l)-r'(v+k)r(u+l)\bigr]\\
 &\quad+\bigl[r'(u+l)r(v+k)-r'(v+l)r(u+k)\bigr]\ge0.
\end{align*}
Each bracket is nonnegative because \(r'\) is nonnegative and decreasing
and \(r\) is increasing. Pairing the \(k\) and \(n-k\) terms in (9)
therefore proves balancing monotonicity for any nonnegative symmetric
weights; the central term, if present, is half such a pair.
Apply this to \(w_k=g_kg_{n-k}\) in (15). Finally, \(d=0\) makes \(G\)
independent of \(\mu\), so (1) vanishes. Replacing \((a,c)\) by
\((a+t,c+t)\) also gives the inequality after any common base shift
\(t\ge0\), as required by the Wright definition. All cases are covered.
\end{proof}

\noindent\textbf{Verification and scope.}
The proof above is analytic for every allowed parameter and coefficient
sequence. The companion program checks 216 independent beta moment
identities, 24 coin-kernel cases, 225 exact Bernstein-density cases,
9 gamma integral evaluations, and 135 full coefficients. These finite
checks support the identities; they are not substitutes for the proof.
The final status and claim checks are recorded separately in the
literature audit.

\paragraph{Task 2.}
Characterize all symmetric nonnegative weight vectors \(w\) for which
(9) increases under balancing for every \(u,v,d>0\), and determine
whether condition (8) can be weakened sharply.

\paragraph{Task 3.}
Develop a basic-hypergeometric counterpart of the beta-mixture argument,
and prove or refute the corresponding real-shift inequalities for
series formed with \(q\)-gamma ratios.

\paragraph{Task 4.}
Determine the exact equality conditions in (1) coefficient by
coefficient, and obtain quantitative lower bounds from the
concentration comparisons used in the proof.

\paragraph{Task 5.}
Extend the concentration argument to Dirichlet mixtures in three or more
coordinates, and identify coefficient conditions that make the resulting
gamma-ratio sums monotone under majorization of their parameters.

\begin{thebibliography}{9}
\bibitem{KK}
S.~I. Kalmykov and D.~B. Karp,
Log-convexity and log-concavity for series in gamma ratios and applications,
\emph{J. Math. Anal. Appl.} \textbf{406} (2013), 400--418.
\href{https://doi.org/10.1016/j.jmaa.2013.04.061}{doi:10.1016/j.jmaa.2013.04.061};
\href{https://arxiv.org/abs/1211.2882}{arXiv:1211.2882}.
\end{thebibliography}
\end{document}
